\section{Conclusion}
\label{sec:conclusion}

Gaussian noise can be small in every resolved direction and still become a
large perturbation as resolution grows. This paper identifies the missing
criterion: dimension-uniform $\Schatten{2}$ control is the sharp threshold for
refinement-stable Gaussian covariance objectives and their sample-level
fluctuations. Schatten-filtered noise makes that requirement explicit, while
the Carleman--Fredholm determinant supplies the correct objective in the
Hilbert--Schmidt but non-trace-class regime. The resulting view explains why
frequency attenuation and dimensionality reduction alone do not determine
continuum stability, and turns the operator-ideal role of latent encoders into
a measurable hypothesis. Noise injection should therefore be designed and
audited as a resolution-indexed spectral object, not as a dimension-free scalar
hyperparameter.
